\chapter{Grammars, states and histories}\label{ch:grammars}

This chapter fixes the objects of the monograph. We define generative systems abstractly, specialise to \emph{commutative realizations} (in which the endpoint depends only on seed and content), present the arithmetic rule SELECT-II and the symbolic event grammars as two realizations, introduce the admissible-set complexes that carry the topology, and state precisely which constructions of the later chapters need integer divisibility and which work for arbitrary finite sets of histories.

\section{Generative systems, histories and fibres}
\begin{definition}[generative system]\label{def:grammar}
A \emph{generative system} $\mathcal G=(\Lambda,X,X_0,(\tau_\ell)_{\ell\in\Lambda})$ consists of a set $\Lambda$ of \emph{labels}, a set $X$ of \emph{states}, a set $X_0\subseteq X$ of \emph{sources}, and for each label a partial map $\tau_\ell\colon X\rightharpoonup X$. A \emph{history} is a pair $h=(x_0;\ell_1,\dots,\ell_n)$ with $x_0\in X_0$ such that $x_i=\tau_{\ell_i}(x_{i-1})$ is defined for $1\le i\le n$; the word $w(h)=(\ell_1,\dots,\ell_n)$ is its \emph{word}, $x_n$ its \emph{endpoint} $t(h)$, and the multiset $\{\ell_1,\dots,\ell_n\}$ its \emph{content} $Q(h)$. The \emph{fibre} of $t\in X$ is the set $F_t$ of histories with endpoint $t$, and $F_t(x_0)$ the subset with source $x_0$. When $F_t$ is finite and nonempty, the \emph{(Hartley) ambiguity} of $t$ is $A(t)=\log_2|F_t|$.
\end{definition}

A fibre is a disjoint union of \emph{pieces} $P(x_0,Q)$ of histories with a given source and content. We write $\Ent(\cdot)$ for Shannon entropy in bits and use the \emph{uniform ensemble} on a finite fibre unless another probability is stated (\cref{sec:probconv}).

\begin{definition}[commutative realization, legal swap]\label{def:swap}
A generative system is a \emph{commutative realization} if $X$ lies in an abelian group (written multiplicatively) and $\tau_\ell(x)=x\,r_\ell$ for a fixed $r_\ell$ on a domain $\mathrm{dom}\,\tau_\ell\subseteq X$. Then $t(h)=x_0\prod_i r_{\ell_i}$ depends only on the source and content. A \emph{legal adjacent swap} exchanges two adjacent distinct labels $\ell_i\ne\ell_{i+1}$ of a history when the exchanged word is still a history from the same source. The \emph{swap graph} $G_t$ of a fibre has the histories as vertices and the legal swaps as edges; its connected components are the \emph{recombination classes}, and the class of $h$ is written $C(h)$.
\end{definition}

\begin{lemma}\label{lem:swap-invariants}
In a commutative realization a legal swap preserves source, content and endpoint. Hence the random variables $T$ (endpoint), $(M,Q)$ (source and content), $C$ (class) and $W$ (the history) satisfy that each determines the previous one: $T=f_1(M,Q)$, $(M,Q)=f_2(C)$, $C=f_3(W)$.
\end{lemma}
\begin{proof}
Swapping two labels permutes the word, so the multiset of labels and the source are unchanged, and $t=x_0\prod r_{\ell_i}$ is unchanged by commutativity of the group. Source and content are therefore constant on connected components of $G_t$, hence functions of the class; the class is a function of the history; the endpoint is a function of source and content.
\end{proof}

\section{The arithmetic realization SELECT-II}\label{sec:selectII}
The first arithmetic realization is the frozen UNNS operator SELECT-II. We use it as the main worked example for two reasons: its histories can be enumerated exactly, and its admissibility is governed by divisibility, which gives a rich family of finite combinatorial structures. We do not suggest that it is the only legitimate generative rule of the UNNS substrate.

\begin{definition}[SELECT-II]\label{def:selectII}
States are positive integers, labels are positive integers $N\ge1$, and
\begin{equation}\label{eq:selectII}
T_N(m)=m\,R(N),\qquad R(N)=\frac{(N+1)^2}{N}=N+2+\frac1N,\qquad \mathrm{dom}\,T_N=\{m: N\mid m\}.
\end{equation}
A source is a positive integer. A history is \emph{complete} if its source is squarefree (we show in \cref{prop:fibre-finite} that this means the source has no predecessor). The \emph{complete-history fibre} $F_t$ consists of all complete histories ending at $t$; the \emph{fixed-source fibre} $F_t(m_0)$ of all histories from $m_0$.
\end{definition}
SELECT-II is a commutative realization with $r_N=R(N)$, so \cref{lem:swap-invariants} applies. Since $R(N)$ is rational and the domain condition is divisibility, a label word is \emph{formally} applicable from any seed but \emph{admissible} only if every intermediate state is an integer divisible by the next label.

\section{Symbolic event grammars}\label{sec:symbolic}
For the topological chapters we do not need an integer state at all. Take a set $[k]=\{0,\dots,k-1\}$ of $k$ \emph{distinguishable events} and consider \emph{event grammars}: sets $A\subseteq S_k$ of admissible orderings. Two families occur:
\begin{itemize}
\item \emph{pairwise grammars}: a digraph $\Gamma$ on $[k]$ and $A=\{$orderings $(x_1,\dots,x_k)$ with $(x_i,x_{i+1})\in\Gamma$ for all $i\}$, the Hamiltonian paths of $\Gamma$; when events carry categories (for instance three categories $S,R,N$, with several events per category) the digraph is the preimage of a rule set on categories;
\item \emph{triple grammars}: $A$ is the set of orderings avoiding a given set of consecutive triples.
\end{itemize}
These grammars have no seed and no arithmetic; they realise the same combinatorial objects as the arithmetic rule (sets of orderings, hence subcomplexes of the permutohedron), and they are used in \cref{ch:locality} to separate what depends on arithmetic from what depends only on the \emph{locality} of the constraints. A natural-language reading (``Sunny, Rainy, Snowy'' with a rule on which weather may follow which) is a symbolic test system only; no weather data enter, and nothing is claimed about meteorology.

\section{Admissible sets and the permutohedral complex}\label{sec:complex}
Fix $k\ge2$. Write $[k]=\{0,\dots,k-1\}$ and let $S_k$ be the set of orderings (permutations) $w=(w_1,\dots,w_k)$ of $[k]$. For $w$ and $1\le j\le k-1$ the \emph{prefix set} is $I_j(w)=\{w_1,\dots,w_j\}$; the chain $I_1(w)\subset\dots\subset I_{k-1}(w)$ is the \emph{flag} of $w$, and $w\mapsto$ flag is a bijection between orderings and maximal chains of the poset $\Pm$ of proper nonempty subsets of $[k]$.

\begin{definition}[ordered partitions, faces]\label{def:faces}
A \emph{face} of the permutohedron $P_k$ is an ordered set partition $F=(B_1|\dots|B_m)$ of $[k]$ into nonempty blocks, $1\le m\le k$; it has $\dim F=k-m$. Its vertex set $V(F)$ is the set of orderings that list $B_1$ first (in any order), then $B_2$, and so on. Faces of dimension $0,1,2,3$ with block sizes $(1,\dots,1)$, $(2)$, $(2,2)$ or $(3)$, and $(2,2,2)$, $(3,2)$ or $(4)$ are vertices, edges, squares or hexagons, and cubes, hexagonal prisms or truncated octahedra. The face $F'$ is a face of $F$ if $F'$ refines $F$. The \emph{proper prefix sets} of $F$ are
\[
\Pref(F)=\{B_{<i}\cup U:\ 1\le i\le m,\ \emptyset\neq U\subseteq B_i\}\setminus\{[k]\},\qquad B_{<i}=B_1\cup\dots\cup B_{i-1}.
\]
\end{definition}
\noindent The 1-skeleton of $P_k$ is the Cayley graph of $S_k$ for the adjacent transpositions $s_1,\dots,s_{k-1}$ (positions $j,j{+}1$); an edge is the face with one block of size two. It is bipartite (by the sign of the permutation). We use the standard facts about permutohedra (\cite{Ziegler}, \cite{BjornerBrenti}): the faces form a regular cell decomposition of the boundary of a convex polytope of dimension $k-1$, and $P_k$ is simple.

\begin{definition}[history complex]\label{def:KA}
For a set $A\subseteq S_k$ of admissible histories, the \emph{history complex} $K_A$ is the full subcomplex of $P_k$ spanned by $A$: a face $F$ of $P_k$ belongs to $K_A$ iff $V(F)\subseteq A$. Its Betti numbers $\beta_j(K_A)=\dim H_j(K_A;\F)$ are the \emph{history Betti numbers}. The 1-skeleton of $K_A$ is the induced subgraph $G_A$ of the permutohedron graph on $A$, the \emph{legal-swap graph}.
\end{definition}
For a commutative realization with \emph{distinct} labels $\ell_1<\dots<\ell_k$ (a \emph{content with distinct labels}; arrangement $=$ ordering of the labels, addressed by rank), the histories of a piece $P(m_0,Q)$ are the admissible orderings $A=A_Q(m_0)\subseteq S_k$, and $G_A$ is the swap graph of the piece because swapping adjacent distinct labels is the transposition of adjacent positions. For contents with repeated labels the swap graph is defined but the permutohedral complex is not (the cell rule needs distinct labels); Betti numbers are then not claimed. All homological statements in this monograph concern distinct-label pieces.

\begin{definition}[prefix-defined admissibility]\label{def:prefix-defined}
A set $A\subseteq S_k$ is \emph{prefix-defined} if there is a family $\mathfrak F\subseteq\Pm$ such that $w\in A$ iff all proper prefix sets $I_1(w),\dots,I_{k-1}(w)$ lie in $\mathfrak F$. A prefix set $I\in\Pm$ is \emph{dead} if every ordering having $I$ as a prefix set (equivalently every flag through $I$) is not in $A$.
\end{definition}

\begin{lemma}[faces of prefix-defined complexes]\label{lem:faces-prefix}
If $A$ is prefix-defined by $\mathfrak F$, then a face $F$ of $P_k$ lies in $K_A$ iff $\Pref(F)\subseteq\mathfrak F$.
\end{lemma}
\begin{proof}
The proper prefix sets of a vertex $w=\pi_1\cdots\pi_m\in V(F)$ ($\pi_i$ an ordering of $B_i$) are the sets $B_{<i}\cup U$ with $U$ a nonempty prefix set of $\pi_i$ (excluding $[k]$). Conversely every nonempty $U\subseteq B_i$ is a prefix set of some ordering of $B_i$ (list $U$ first), and the other blocks may be ordered arbitrarily. Hence $\bigcup_{w\in V(F)}\{I_j(w)\}=\Pref(F)$, and $V(F)\subseteq A$ iff $\Pref(F)\subseteq\mathfrak F$.
\end{proof}

\section{Which constructions need arithmetic?}\label{sec:needs}
\Cref{tab:needs} lists the main constructions, the structure each requires, and where it is used. The table is the contract that later chapters keep: a result in a row is claimed only for the stated class of systems.

\begin{table}[ht]\centering\small
\begin{tabularx}{\textwidth}{@{}>{\raggedright\arraybackslash}p{4.2cm}>{\raggedright\arraybackslash}X>{\raggedright\arraybackslash}p{2.2cm}@{}}
\toprule
Construction & Required structure & Where\\\midrule
Reverse edges, finite fibres, squarefree complete sources & SELECT-II (integers; strict growth $\ge4$ per step) & \cref{ch:ancestry}\\
Ambiguity bounds; chain-rule hierarchy $T\!\subset\!(M,Q)\!\subset\!C\!\subset\!W$ & any finite fibre; the nesting needs a \emph{commutative} realization & \cref{ch:ancestry}\\
Swap graph, accessibility profile, radius-two identity & any finite set of orderings with adjacent-transposition swaps & \cref{ch:swaps,ch:access}\\
Lemma on prefix sets (admissible iff denominators divide) & SELECT-II (divisibility); the \emph{consequence} (prefix-defined) is all that later chapters use & \cref{ch:arith}\\
Block factorization, cube-product property & any \emph{prefix-defined} $A$ & \cref{ch:arith,ch:swaps}\\
lcm closure, minimal seed, divisibility filtration, survival laws, gates & SELECT-II (divisibility lattice) & \cref{ch:arith,ch:topology,ch:atlas}\\
Reconstruction of $K_A$ from the coloured graph & arbitrary $A\subseteq S_k$ & \cref{ch:swaps}\\
Two-skeleton from the uncoloured graph & cube-product property (all prefix-defined $A$), $k\le8$ & \cref{ch:swaps}\\
Lift operator and its boundary identity & arbitrary chains on permutohedra (pure combinatorics); \emph{survival} needs arithmetic & \cref{ch:dimension}\\
Alexander duality $K_A\leftrightarrow Y_A$ & arbitrary $A\subsetneq S_k$ nonempty & \cref{ch:duality}\\
Dead-set nerve theorem & prefix-defined $A$, any $k$ for the homotopy statement; $\beta_2$ formula for $k=6$ & \cref{ch:duality}\\
Pairwise locality theorem & Hamiltonian paths of an arbitrary digraph on $k=6$ events (no arithmetic) & \cref{ch:locality}\\
\bottomrule
\end{tabularx}
\caption{Constructions and what they require. ``Prefix-defined'' includes SELECT-II (\cref{lem:F}) but is a much larger class.}\label{tab:needs}
\end{table}

\section{Probability conventions}\label{sec:probconv}
Three different probabilities occur and are never mixed.
\begin{enumerate}[label=(\roman*)]
\item \emph{Ensemble on a fibre.} Unless stated otherwise, a fibre $F_t$ carries the uniform distribution; then $\Ent(W|T=t)=A(t)$ and every entropy in Part~I is an exact finite sum of terms $\frac{n}{N}\log_2\frac Nn$ with integer $n,N$. A nonuniform distribution is a modelling choice and is stated each time it occurs; it never exceeds $A(t)$ (\cref{prop:entropy-bounds}).
\item \emph{Ensemble over systems.} Part~IV treats a \emph{system} (a content together with a seed regime) as a random variable $X$ and asks how much an observation $O(X)$ tells about a property $Y(X)$. The \emph{natural} ensemble is uniform over the systems of a stated census; the \emph{balanced} ensemble gives equal total weight to $Y=1$ and $Y=0$. These are experimental conventions, not probability laws of the grammar.
\item \emph{Generalisation.} Entropies computed on the whole finite ensemble by exact class counts are \emph{in-sample} (plug-in); they measure collisions of signature classes. \emph{Out-of-sample} information is the information gained by a rule trained on one sample and tested on another. The two are different and are compared in \cref{ch:predictive}.
\end{enumerate}
The Hartley quantity $A(t)=\log_2|F_t|$ needs no probability; it equals the Shannon entropy of the uniform ensemble and is the maximum of $\Ent(W|T=t)$ over all distributions on $F_t$.
